\documentclass[10pt,a4paper]{amsart}
\usepackage[margin=19mm]{geometry}
\usepackage{amsmath,amssymb,mathtools}
\usepackage[scr=rsfs,cal=euler]{mathalfa}
\usepackage{xfrac}
\usepackage[unicode,hidelinks,bookmarks=false]{hyperref}
\newcommand{\ie}{i.e.\ }
\newcommand{\cf}{cf.\ }
\newcommand{\resp}{respectively\ }
\newcommand{\Subsection}{\subsection}
\providecommand{\nobreakdash}{\nobreak\hbox{-}}
\setlength{\emergencystretch}{2em}
\multlinegap0pt
\newtheorem{lem}{Lemma}
\title{Extremal representations of functions of matrices and applications to multivariate prediction}
\author{Michael Frank, Lutz Klotz, Andreas Lasarow}
\date{Source: arXiv:2606.19359v1 (CC BY 4.0)}
\begin{document}
\maketitle

\noindent\emph{Source and licence.} Extremal representations of functions of matrices and applications to multivariate prediction. Version arXiv:2606.19359v1 (CC BY 4.0), distributed by arXiv under the Creative Commons Attribution 4.0 International licence, http://creativecommons.org/licenses/by/4.0/, as recorded in the arXiv licence metadata for this exact version and shown on the abstract page https://arxiv.org/abs/2606.19359v1. Full licence notice: \texttt{licenses/arxiv-2606.19359-CC-BY-4.0.txt}.\par\smallskip

\noindent\textit{Editorial scope.} Counted unit: the article's lem l5.8 (source0.tex lines 939--948). The article's complete original proof of it is delivered with it (source0.tex lines 950--981, one \texttt{proof} environment of 32 lines). No mathematics is written, repaired or re-derived: the statement and the proof are the article's own lines, in the article's own notation, and no retained line is substituted or re-flowed.  Retained uncounted as the source-local context the proof needs: lines 934--936.  Nothing was omitted from any retained span, so the package carries no omission record.  No retained line contains a citation, so the wrapper carries no bibliography and no bibliography entry.  No retained line contains a figure environment, an image-inclusion command or drawing code, so no image is delivered and the package declares no asset dependency.  Editorial formatting only, in addition: an AMS \texttt{amsart} preamble that restates the article's own declarations and macros used by the retained lines, namely the environments \texttt{lem} on separate counters, together with the standard packages those lines need.  The retained environments print in this document as Lemma 1 in order of appearance, whereas the article prints the same items with the numbering of its own full document.  The complete native source of the article as distributed in the arXiv e-print for this exact version is bundled unchanged as \texttt{sources/203138/source0.tex}.
\begin{equation} \label{f5.16}
     {\rm Find} \,\, \mu_n = \min f_n(x_1,\ldots,x_n) \,\, {\rm under} \, {\rm the} \, {\rm condition} \,\, E_{2,n}(x_1,\ldots,x_n)=1 \, .
\end{equation}     

\begin{lem} \label{l5.8}
   Let $n \geq 3$. The problem ${\mathcal E}_n$ has at most one stationary point $(x_1^{(n)},\ldots,x_n^{(n)} )$. The stationary point exists if and only if 
   \begin{equation}   \label{f5.17}
       \min \left\{ \sum_{j=1}^n a_j - (n-1) a_n , \left(\sum_{j=1}^n a_j \right)^2  - (n-1)\sum_{j=1}^n a_j^2 \right\} > 0 \, .
   \end{equation}
In this case
   \begin{equation}   \label{f5.18}
        \mu_n= f(x_1^{(n)},\ldots,x_n^{(n)}) = \sqrt{\frac{2}{n-1}} \left( \left(\sum_{j=1}^n a_j \right)^2 - (n-1) \sum_{j=1}^n a_j^2 \right)^{1/2} \, .
   \end{equation}     
\end{lem}

\begin{proof}
Let $F(x_1,\ldots,x_n,\lambda) = f(x_1,\ldots,x_n) - \lambda (E_{2,n}(x_1,\ldots,x_n)-1)$, add the $n$ equations $\frac{{\partial} F}{{\partial} x_k} (x_1,\ldots,x_n,\lambda) = a_k-\lambda (\sum_{j=1}^n x_j - x_k)=0$, $k \in \{ 1,\ldots,n \}$, and subtract from this sum the $k$-th equality multiplied by $(n-1)$. One obtains $\sum_{j=1}^n a_j - (n-1)a_k = (n-1) \lambda x_k$ what transforms to
\begin{equation}   \label{f5.19}
    x_k = ((n-1)\lambda)^{-1} \left(\sum_{j=1}^n a_j - (n-1)a_k \right) \, .
\end{equation}
Since $x_n$ should be positive, necessarily $\sum_{j=1}^n a_j - (n-1)a_n \not= 0$. Moreover, if this expression would be negative then necessarily $\lambda < 0$, leading to an contradiction $x_1 <0$ since  $\sum_{j=1}^n a_j - (n-1)a_n > 0$. Consequently, the condition 
$\sum_{j=1}^n a_j - (n-1)a_n  > 0$ is necessary for the existence of a stationary point.  Using the equality $E_{2,n} (x_1,\ldots,x_n) = \frac{1}{2} \, \left((\sum_{j=1}^n x_j)^2 - (n-1) \sum_{j=1}^n x_j^2 \right)$ one easily obtains
\begin{equation}   \label{f5.20_1}
    \lambda^2 =(2 (n-1)^2)^{-1} \left((\sum_{j=1}^n a_j)^2 - \sum_{j=1}^n a_j^2 \right) \, ,
\end{equation}
by \eqref{f5.16} and \eqref{f5.19}, what implies inequality $\left((\sum_{j=1}^n a_j)^2 - (n-1) \sum_{j=1}^n a_j^2\right) > 0$ to be also necessary for the existence of a stationary point.

Conversely, if \eqref{f5.17} is satisfied, the computations above, particularly \eqref{f5.20_1} and \eqref{f5.19}, reveal the existence of a unique stationary point $(x_1^{(n)},\ldots,x_n^{(n)})^\top$. One can compute $f_n(x_1^{(n)},\ldots,x_n^{(n)}) = \sqrt{2/(n-1)} \left((\sum_{j=1}^n a_j)^2 - (n-1) \sum_{j=1}^n a_j^2 \right)^{1/2}$. We have only to demonstrate finally that this term is the minimum $\mu_n$ of the extremal problem ${\mathcal E}_n$. Since, obviously, the minimum of ${\mathcal E}_n$ at the boundary $\{ (x_1,\ldots,x_n)^\top : x_j=0 \,\, {\rm for} \, {\rm some} \,\, j \in \{1,\ldots,n)\}\}$ equals the minimum at the set $\{ (x_1,\ldots,x_{n-1},0): (x_1,\ldots,x_{n-1})^\top \in [0,+\infty)^{n-1} \}$ it suffices to show that for all $n \geq 3$ the following fact is true: If ${\mathcal E}_n$ has a stationary point $(x_1^{(n)}, \ldots , x_n^{(n)} )$ then ${\mathcal E}_{n-1}$ has a stationary point $(x_1^{(n-1)}, \ldots , x_{n-1}^{(n-1)} )$ as well, and 
\begin{equation}  \label{f5.20_2}
    f_n(x_1^{(n)},\ldots,x_n^{(n)}) \leq f_{n-1}(x_1^{(n-1)},\ldots,x_{n-1}^{(n-1)}) \, .
\end{equation}    
We proceed by induction on $n$: If ${\mathcal E}_3$ has a stationary point $(x_1^{(3)},x_2^{(3)},x_3^{(3)} \}$ then 
\begin{eqnarray*}
          f_3(x_1^{(3)},x_2^{(3)},x_3^{(3)}) & = & ((a_1+a_2+a_3)^2 -2(a_1^2+a_2^2+a_3^2))^{1/2} \\
                                                               & = &  (4a_1a_2 - (a_1+a_2-a_3)^2)^{1/2} \\
                                                            & \leq & 2 \sqrt{a_1a_2} = f_2(x_1^{(2)},x_2^{(2)}) \, .
\end{eqnarray*}
For arbitrary $n$, the inequality $a_{n-1} \leq a_n$ yields
\begin{equation} \label{f5.21}
      \sum_{j=1}^n a_j - (n-1) a_n \leq \sum_{j=1}^{n-1} a_j - (n-2) a_{n-1} \, .
\end{equation}
Moreover, the valid inequality $0 \leq (\sum_{j=1}^{n-1} a_j - (n-2) a_n)^2$ can be transformed into the equivalent inequality
\begin{equation}    \label{f5.22}
      ( \sum_{j=1}^n a_j )^2 - (n-1) \sum_{j=1}^n a_j^2 \leq \frac{n-1}{n-2} \left(  (\sum_{j=1}^{n-1} a_j)^2 - (n-2) \sum_{j=1}^{n-1} a_j^2\right) \, .
\end{equation}
Therefore, by \eqref{f5.18}, \eqref{f5.21} and \eqref{f5.22} the existence of a stationary point of ${\mathcal E}_n$ is sufficient for the existence of a stationary point of ${\mathcal E}_{n-1}$. Finally, by \eqref{f5.19} and \eqref{f5.22} the inequality \eqref{f5.20_2} follows. 
\end{proof}   

\end{document}
